% Source: https://artofproblemsolving.com/community/c3h1577812_middle_school_math
% Stable identity: aops-topic-1577812-q01
Let $S_n=1+\dfrac{1}{2}+\dfrac{1}{4}+\cdots+\dfrac{1}{2^{n}}$ be the sum of the first $n+1$ terms. Then
\[
\dfrac{1}{2}S_n=\dfrac{1}{2}+\dfrac{1}{4}+\cdots+\dfrac{1}{2^{n+1}},\qquad\text{so}\qquad S_n-\dfrac{1}{2}S_n=1-\dfrac{1}{2^{n+1}},
\]
that is, $S_n=2-\dfrac{1}{2^{n}}$. Every partial sum falls short of $2$ by exactly the last term added, and this shortfall tends to $0$. Hence the infinite sum (the limit of the partial sums) is
\[
\boxed{2}.
\]
(This is the geometric series with first term $1$ and ratio $\dfrac{1}{2}$, whose sum is $\dfrac{a}{1-r}=\dfrac{1}{1-\dfrac{1}{2}}=2$ with $a=1$, $r=\dfrac{1}{2}$.)
